% Generated draft: requires mathematical and source review.
\documentclass[11pt]{article}
\usepackage[margin=25mm]{geometry}
\usepackage{fontspec}
\setmainfont{Latin Modern Roman}
\usepackage{amsmath,amssymb,mathtools,unicode-math}
\setmathfont{Latin Modern Math}
\usepackage{enumitem,needspace,xurl,hyperref}
\setlength{\parindent}{0pt}
\setlength{\parskip}{5pt}
\setlength{\emergencystretch}{3em}
\newcommand{\N}{\mathbb{N}}
\newcommand{\Z}{\mathbb{Z}}
\newcommand{\Q}{\mathbb{Q}}
\newcommand{\R}{\mathbb{R}}
\newcommand{\C}{\mathbb{C}}
\newcommand{\F}{\mathbb{F}}
\newcommand{\A}{\mathbb{A}}
\newcommand{\PP}{\mathbb{P}}
\newcommand{\E}{\mathbb{E}}
\newcommand{\eps}{\varepsilon}
\newcommand{\dd}{\,\mathrm{d}}
\begin{document}
\section*{UnsolvedMath ID: 10400108}
\textbf{Source problem code:} catalogue entry
\subsection{Definitions and mathematical statement}
For a prime \(p\), let the dihedral rack/quandle \(R_p\) be \(\mathbb Z/p\mathbb Z\) with \(x*y=2y-x\). In the rack-homology interpretation supplied by the summary, take \(C_n^R(R_p)=\mathbb Z[R_p^n]\) and boundary
\[
 \partial_n(x_1,\ldots,x_n)=\sum_{i=2}^n(-1)^i
 \left[(x_1,\ldots,\widehat{x_i},\ldots,x_n)
 -(x_1*x_i,\ldots,x_{i-1}*x_i,x_{i+1},\ldots,x_n)\right].
\]
The degree-three group is \(H_3^R=\ker\partial_3/\operatorname{im}\partial_4\), with integer coefficients.

The extracted conjecture says
\[
 H_3(R_p)\cong\mathbb Z\oplus\mathbb Z/p\mathbb Z
 \quad\text{for }p\text{ prime}.
\]
Its homology superscript is missing, while the summary identifies rack homology and reports the formula proved for odd primes. The prime \(2\) cannot be silently included under the displayed rack convention: \(R_2\) has trivial operation, all boundaries vanish, and \(H_3^R(R_2)\cong\mathbb Z^8\). Thus the all-prime wording needs correction or a different intended convention; the coherent reported theorem is the odd-prime rack statement.

The summary attributes that odd-prime result to Niebrzydowski--Przytycki, using Mochizuki's mod-\(p\) quandle computation and a rack/quandle homology decomposition. It also reports \(p\)-torsion in higher rack degrees \(n\ge3\) and related \(R_3\) results, without exact subsidiary hypotheses. Those auxiliary reports are not substituted for the degree-three target or independently certified. The registry solved label is retained as a literature attribution, with the exceptional-prime issue explicit.
\subsection{Short English statement}
Compute the third integer rack homology of the dihedral quandle. The supplied literature report proves \(\mathbb Z\oplus\mathbb Z/p\) for odd primes; the extracted all-prime wording fails at \(p=2\) under the standard rack definition.
\subsection{Sources}
\begin{itemize}
\item Local registry: \url{data/unsolvedmath/problems.json}, permanent UnsolvedMath ID \(10400108\), source code catalogue entry; supplied statement, background, and research summary, updated 2026-08-06.
\item Ohtsuki, \emph{Problems on invariants of knots and 3-manifolds} (2002), Conjecture 5.12, attributed to R. Fenn, C. Rourke and B. Sanderson, PDF page 93: \url{https://msp.org/gtm/2002/04/gtm-2002-04-024s.pdf}. Supplied PDF-font-structure extraction accessed 2026-07-29.
\item M. Niebrzydowski and J. H. Przytycki, ``Homology of dihedral quandles'', \emph{Journal of Pure and Applied Algebra} 213 (2009), 742--755; arXiv:math/0611803 (2006), as supplied by the summary.
\item Mochizuki's computation \(H_3^Q(R_p;\mathbb Z_p)\cong\mathbb Z_p\) and a rack/quandle splitting are cited by the summary without separate full references.
\end{itemize}
\end{document}
